\documentclass[11pt]{article}
\usepackage[margin=1in]{geometry}
\usepackage{amsmath,amssymb,booktabs,hyperref,natbib}
\usepackage[T1]{fontenc}

% =================================================================
% GAPS.tex -- Validation notes on the Problem Formulation.
%
% Supersedes GAPS.md (kept as markdown was the wrong format for a document
% whose whole content is thesis equations, cross-references, and a table
% meant to be pasted into the thesis's own appendix -- it belongs typeset
% alongside the thesis it discusses, not as a loose Markdown file).
%
% This file stands alone (compiles with `pdflatex GAPS.tex`) but its
% equation numbers/labels intentionally mirror problem_formulation_fixed.tex,
% so \eqref calls below refer to that file's labels, not local equations,
% wherever they cite eq:xxx from Section IV. Re-numbering is not attempted
% here; this is a companion document, not a subfile of the thesis itself.
%
% Three parts:
%   Part I   -- internal-consistency review of problem_formulation.tex
%              ([G1]-[G17], matching the inline tags in the fixed .tex)
%   Part II  -- comparison against Knippenberg, Holenderski & Menkovski (2021),
%              the closest prior-art paper this section's environment and
%              observation model is adapted from ([K1]-[K8])
%   Part III -- comparison against Ma, Li, Kumar & Koenig (2017), the paper
%              this section's MAPD terminology and well-formedness condition
%              are adapted from ([M1]-[M4])
%   Appendix -- the notation-disambiguation table, appendix-ready LaTeX to
%              paste into the thesis's own appendix chapter (not present in
%              this repo) rather than left inline in Section IV, per the
%              Problem Formulation / Method / Implementation contract in
%              Method_and_Implementation.md.
% =================================================================

\title{Validation Notes on the Problem Formulation}
\author{}
\date{}

\begin{document}
\maketitle

\tableofcontents

\section*{How to use this document}
\texttt{problem\_formulation\_fixed.tex} is a complete, label-preserving
drop-in revision of \texttt{problem\_formulation.tex}. Every fix there is
tagged inline (\texttt{\% [Gxx]}, \texttt{\% [Kxx]} or \texttt{\% [Mxx]});
this document is the rationale for each tag, split into three parts by where
the finding came from: Part~\ref{sec:internal} was found by checking the
section against itself (does every symbol get defined, do the claimed
properties actually follow); Part~\ref{sec:knippenberg} was found by checking
it against the paper its environment and observation model is adapted from;
Part~\ref{sec:ma2017} was found by checking it against the paper its MAPD
terminology and well-formedness condition are adapted from. Severity key used
throughout: \textbf{Error} = mathematically wrong as stated; \textbf{Gap} =
not wrong, but undefined/underspecified where the section otherwise commits
to full formality; \textbf{Deviation} = a considered, defensible difference
from a source, not necessarily a defect; \textbf{Clarity} = correct and
complete, but hard to read.

\part*{}
\section{Internal-consistency review}
\label{sec:internal}

\subsection*{[Error] The potential-function claim (target-distance annotation)}
\textbf{Original text:} ``The annotation is a potential: its difference along
an edge, $\eta_i(w,t) - \eta_i(v,t)$, is negative exactly on shortest-path
moves.''

\textbf{Why it's wrong:} by the triangle inequality on $d_G$,
\[
  \eta_i(w,t) - \eta_i(v,t) \;\ge\; -c(v,w) \qquad \text{for every } (v,w)\in E,
\]
with \emph{equality iff} $(v,w)$ lies on a shortest path from $v$ to
$q_i(t)$. ``Negative'' is necessary for a shortest-path move but not
sufficient: a move can strictly satisfy the inequality (decrease $\eta_i$ by
less than $c(v,w)$) while not lying on any shortest path. Concretely:
$c(v,w)=1$, $d_G(v,q)=5$, and some other route makes $d_G(w,q)=4.5$ -- the
difference is $-0.5$ (negative), but $(v,w)$ is not on a shortest path (that
would require $d_G(w,q)=4$ exactly).

\textbf{Fix:} restated as the inequality with the correct equality condition
(\texttt{eq:potentialbound} in the fixed .tex); surrounding prose now says a
negative difference indicates \emph{progress}, not \emph{optimality}. This is
the one place a formally stated property was actually false, not merely
underspecified.

\subsection*{[G1] \texorpdfstring{$V_{mov}$}{Vmov} vs. \texorpdfstring{$V_{str}$}{Vstr}/\texorpdfstring{$V_{del}$}{Vdel} -- can an agent ever reach a pickup or delivery vertex?}
Agents only ever occupy $\ell_i(t) \in V_{mov}$. The original text declares
$V_{str}, V_{del} \subseteq V$ (not $\subseteq V_{mov}$) while calling all
four roles a ``partition'' -- which $V_{ep}\subseteq V_{mov}$ already
contradicts, since a partition requires disjoint blocks. If $V_{str}$ and
$V_{del}$ are disjoint from $V_{mov}$, as ``partition'' suggests, an agent can
never occupy a pickup or delivery vertex, leaving pickup/delivery physically
undefined. \textbf{Fix:} $V_{str}, V_{del} \subseteq V_{mov}$ now stated
explicitly (they must be occupiable), ``partition'' replaced by ``label,''
and pairwise disjointness explicitly not assumed.

\subsection*{[G2] The task tuple is missing the SKU it names}
$\tau_j = (r_j, s_j, g_j)$ in the original, but the coupling condition
immediately after it, $x_{r_j}(k_j, s_j) > 0$, uses $k_j$ -- the requested
SKU -- which is never introduced anywhere. The thesis's central equation
rests on an undefined symbol. \textbf{Fix:} \texttt{eq:task} is now a
quadruple $\tau_j = (r_j, s_j, g_j, k_j)$, $k_j \in \mathcal{K}$.

\subsection*{[G3] \texorpdfstring{$\Pi$}{Pi} and \texorpdfstring{$\mathcal{F}$}{F} are never defined as sets}
The formal decision problem optimises over $F \in \mathcal{F}$, $\pi \in
\Pi$, but neither space's domain/codomain is ever written down -- only
\emph{what information} each may condition on (Table~I) is described.
\textbf{Fix:} $\mathcal{F}$ defined as the set of measurable functions
$F : \mathcal{X} \times \mathbb{R}_{\ge0}^{|\mathcal{K}|} \times
\mathbb{R}_{\ge0}^{|E|} \times \mathbb{R}_{\ge0}^{|E|} \to \mathcal{X}$;
$\Pi$ defined as the set of pairs $\pi=(\pi^{\text{assign}},
\pi^{\text{route}})$ of an assignment map and a routing map with stated
domains, constrained only by Table~I. A related loose end fixed in the same
place: eq.~8 ($x_t \in \arg\min$) can be set-valued while eq.~7 ($x_t =
F(\ldots)$) claims a function; the fix adds a tie-breaking-rule requirement
so eq.~7 is well-defined.

\subsection*{[G4] No initial storage configuration}
$x_t = F(x_{t-\Delta}, \ldots)$ needs $x_0$ to bootstrap it at $t=\Delta$;
$x_0$ is never introduced. \textbf{Fix:} $x_0 \in \mathcal{X}$ stated as
supplied with the problem instance.

\subsection*{[G5] \texorpdfstring{$\mathcal{Q}_t, \mathcal{B}_t, \mathcal{C}_t$}{Qt, Bt, Ct} are only described narratively}
Disjointness is asserted (``a property asserted at runtime'') rather than
derived, unlike every other object in the section. \textbf{Fix:} introduces
the assignment time $y_j$ ($r_j \le y_j \le d_j$) and defines
\[
\mathcal{Q}_t = \{\tau_j : r_j \le t < y_j\}, \quad
\mathcal{B}_t = \{\tau_j : y_j \le t < d_j\}, \quad
\mathcal{C}_t = \{\tau_j : d_j \le t\}.
\]
Disjointness now follows from the three intervals partitioning
$[r_j,\infty)$, rather than being separately asserted.

\subsection*{[G6] The congestion feature \texorpdfstring{$\delta_i(t)$}{delta\_i(t)} is defined but never wired in}
$\delta_i(t)$ appears in the Learning Problem subsection, but $o_i(t)$ was
fixed earlier and never amended, and $\delta_i(t)$ is never connected to
$R_i$ either. \textbf{Fix:} moved \texttt{eq:congestion} into the
Observations subsection, right after $\eta_i$, and added $\delta_i(t)$ as a
third component of \texttt{eq:observation}.

\subsection*{[G7] \texorpdfstring{$W_T$}{WT} has no formula and an ambiguous relationship to \texorpdfstring{$\hat{w}_t(e)$}{what(e)}}
$W_T$ appears only inside the system-level criterion, described as having an
operational definition deferred to Implementation, with zero conceptual
description here -- and a name close enough to the already-used per-edge
estimate $\hat{w}_t(e)$ that the two could be mistaken for the same object.

\textbf{First-pass fix (superseded below):} an initial revision gave
$W_T = \frac{1}{|\mathcal{C}_T|}\sum_{\tau_j \in \mathcal{C}_T}(y_j - r_j)$,
mean pre-assignment queueing delay, using $y_j$ from [G5]. This was internally
consistent but, once \texttt{Method\_and\_Implementation.md}'s operational
suggestion for $W_T$ (``timesteps in which an agent waited or was overridden
while holding an active task'') was checked against it, turned out to describe
a \emph{different} quantity -- and queueing delay is largely redundant with the
backlog term $B_T = |\mathcal{Q}_T|+|\mathcal{B}_T|$ already in
\texttt{eq:functional}, the same double-counting concern the section already
raises for throughput vs.\ $\bar\zeta_T$.

\textbf{Current fix:} $W_T$ now measures time lost \emph{after} assignment
instead, which is not redundant with $B_T$: for agent $i$ serving $\tau_j$ over
$t\in[y_j,d_j)$, let $\omega_j(t)\in\{0,1\}$ indicate agent $i$ is blocked at
$t$ (chose wait, or had a move overridden into a wait by the conflict-resolution
operator, Section~\ref{sec:method-mapping}/\texttt{sec:method:resolution}).
Then
\[
W_T = \frac{1}{|\mathcal{C}_T|} \sum_{\tau_j\in\mathcal{C}_T}
      \sum_{t=y_j}^{d_j-1} \omega_j(t).
\]
This version also gives the resolution operator's ``overridden'' flag (needed
anyway for the feasibility mechanism) a direct second use as $W_T$'s data
source, rather than requiring separate instrumentation. $W_T$ remains distinct
from $\hat w_t(e)$ for the same reason as before: one is a realised scoring
statistic, the other a forward-looking estimate feeding $F$.

\subsection*{[G8] Entropy is undefined at the boundary}
\texttt{eq:entropy} is stated ``for $|E_T^+| > 1$'' only; at
$|E_T^+|\in\{0,1\}$ the denominator $\log|E_T^+|=0$. \textbf{Fix:}
$|E_T^+|=1$ defined as maximally concentrated by convention ($H_T := 0$,
$C_T := 1$); $|E_T^+|=0$ reported separately, like $|\mathcal{C}_T|=0$.

\subsection*{[G9] Are constraints (4)--(5) actually constraining anything in the decision problem?}
Section~B already states vertex/swap-conflict feasibility is enforced by a
mechanism applied regardless of which $\pi$ is chosen. If so, every $\pi$
satisfies them by construction, and listing them in the $\arg\min$ over
$(F,\pi)$ doesn't prune the search space. \textbf{Fix:} one clarifying
paragraph: these constraints are retained to state explicitly that a
solution is a claim about \emph{feasible realised trajectories}, unlike
\texttt{eq:feasiblestorage} and \texttt{eq:coupling}, which genuinely do
restrict $F$ and the realised task stream.

\subsection*{[G10] Is \texorpdfstring{$G$}{G} directed, and is it connected?}
One-way aisles imply $G$ must be directed, but this is never stated; $d_G$ is
used pervasively as if always finite, with no connectivity assumption
stated. \textbf{Fix:} $G$ explicitly stated as directed; a finiteness
assumption added (tightened further in Part~\ref{sec:knippenberg}, [K1]).

\subsection*{[Gap] \texorpdfstring{$\mathcal{U}_i$}{Ui} (POSG action space) vs. \texorpdfstring{$\mathcal{U}(v)$}{U(v)} (vertex-local action set)}
\texttt{eq:actions} is explicitly vertex-dependent, but the POSG tuple lists
a fixed per-agent $\mathcal{U}_i$, with no statement of how one becomes the
other. \textbf{Fix:} $\mathcal{U}_i := \bigcup_{v\in V_{mov}}\mathcal{U}(v)$,
with only $\mathcal{U}(\ell_i(t))$ enabled at a given state (the rest masked
to probability zero).

\subsection*{[Gap] Does \texorpdfstring{$J(\theta)$}{J(theta)} apply to all three controller architectures?}
$\theta$ is only introduced for the decentralised architecture, but
\texttt{eq:objective} is presented as though general. \textbf{Fix:} one
sentence: it applies to whichever regime is realised as a learned
controller, with $\theta$ ranging over that regime's own parameters.

\subsection*{[G13] Is the conflict-resolution mechanism itself deterministic?}
Section~J claims execution is deterministic with $P$ stochastic only through
order arrivals, but never confirms the yield-resolution mechanism is itself
deterministic -- if it used randomised tie-breaking, that would be a second
stochasticity source. \textbf{Fix:} one sentence confirming it is
deterministic given the state.

\subsection*{[G11 / Clarity] Symbol collisions: five unrelated concepts spelled with the letter ``C,'' three with ``Q,'' two with ``O''}
Not an error, but slows careful reading (eq.~21 has both $\bar c_T$ and
$\mathcal{C}_T$ in the same line). \textbf{Fix:} rather than renaming
established symbols this late -- Method/Implementation may already reference
them, and a rename should touch the whole thesis at once, not one section in
isolation -- a disambiguation table is provided in
Appendix~\ref{app:notation} of \emph{this} document, ready to paste into the
thesis's own appendix chapter, with only a one-line pointer left inline in
Section~IV (per \texttt{Method\_and\_Implementation.md}'s own rule that a
full notation table is appendix material).

\subsection*{Left open / not changed}
\begin{itemize}
  \item \textbf{$m$ used before being defined} (Section~A uses it,
    Section~B defines $m=|\mathcal{A}|$) -- fixed with a one-line forward
    pointer rather than reordering subsections.
  \item \textbf{A full rename of the C/Q/O-family symbols} -- flagged via
    the appendix table rather than executed; out of scope for a
    Problem-Formulation-only fix.
  \item \textbf{Unit/scale normalisation across the five terms of
    $\mathcal{J}(F,\pi)$} -- summing time, cost, task-count and
    dimensionless quantities under weights $q_\bullet$ with no stated
    calibration scheme is a Method-level concern (how $q_\bullet$ are
    chosen), not a Problem-Formulation error, but worth having an answer
    ready for: linear scalarisation across unlike units needs either
    normalisation or explicit justification for the weights.
  \item \textbf{Section~H's $u_t\sim\pi(\cdot\mid s_t)$ vs.\ the
    decentralised $\pi_\theta(\cdot\mid o_i(t))$} -- already reconciled
    informally (``with the corresponding product of local policies under
    decentralised execution''); not tightened further.
\end{itemize}

\subsection{[G14] Congestion sensitivity in routing was decentralised-only, which breaks the cross-architecture comparison}
The Project Description added to this file (Section~\ref{sec:pd-aim-and-scope})
lists ``congestion sensitivity: disabled or enabled in the storage and routing
layers'' as an independent variable, orthogonal to controller architecture
(Table~\ref{tab:information}). Storage already had this: $\beta$ in
\texttt{eq:storageobjective} toggles it. Routing did not: $\delta_i(t)$
\texttt{eq:congestion} was defined only for the decentralised agent's local
window, so ``congestion sensitivity: enabled'' had no realisable meaning for
the centralised or section-based architectures at all -- two of the three
levels of the controller-architecture factor could not actually take both
levels of the congestion-sensitivity factor. This is not a readability gap;
it is a factorial-design defect: RQ2 (does the effect of coupling differ by
architecture) and RQ3 (does congestion sensitivity help) cannot both be asked
cleanly if one factor is only realisable at one level of the other.

\textbf{Fix:} \texttt{eq:congestion}'s formula and label are unchanged (so
every existing \eqref{eq:congestion} in \texttt{method\_implementation\_-
skeleton.tex}, i.e.\ \texttt{eq:reward} and \texttt{eq:readout}, still resolves
identically). A general occupancy fraction $\delta(S,t;\mathrm{excl})$,
\texttt{eq:occupancy}, is introduced in \ref{sec:pf:controllers} -- the
earliest point in the document where all three architectures and the zone
partition \texttt{eq:partition} are already in scope -- and $\delta_i(t)$ is
restated as its decentralised instance ($S=V^{(r_{\mathrm{cng}})}_i(t)$,
$\mathrm{excl}=\{a_i\}$), with the centralised and section-based instances
given in prose ($S=V_{\mathrm{mov}}$ and $S=Y_q$ respectively, both with
$\mathrm{excl}=\emptyset$, since a coordinator is not itself an occupant the
way an agent asking about its own neighbourhood is). Concretely realising
this signal inside the centralised/section-based planners --- a
congestion-weighted variant of each, mirroring the storage layer's
$F_{\mathrm{dem}}/F_{\mathrm{cng}}$ ladder --- is left as a named, forward-
pointed Method-level task (\ref{sec:method:controllers}), not written here,
per this document's own Problem-Formulation/Method boundary.

\subsection{[G15] Computational cost was promised in the new Introduction but never formalised}
The Introduction and Project Description name ``computational cost'' /
``controller runtime'' as an outcome measure and dependent variable, but
\texttt{eq:functional} had no term of that kind and \ref{sec:pf:measures}
never defined one. \textbf{Fix:} $\kappa_T$, \texttt{eq:runtime}, mean
per-decision wall-clock cost, defined and explicitly excluded from
\texttt{eq:functional} for the same reason throughput already is (different
unit family; a $q_\bullet$ weight would fix an arbitrary seconds-to-quality
exchange rate). No Method-file edit was needed for this one:
\texttt{method\_implementation\_skeleton.tex}'s \ref{sec:impl:cost} already
specifies wall-clock logging split by $\pi^{\mathrm{assign}}$ vs
$\pi^{\mathrm{route}}$ and training vs inference time -- the operational
``how'' was already in the right place, waiting for a ``what'' to point at it.
\texttt{eq:rqformal} is also extended to range over the congestion-sensitivity
toggle from [G14], closing the RQ3 gap the same review surfaced.

\subsection*{Correction to prior advice in this review: do not trim \ref{sec:pf:observations}}
An earlier round of this discussion (not in this document) suggested
compressing the observation subsection's proof-level detail --
\texttt{eq:potentialbound}'s equality condition, the congestion radius, the
communication graph -- to rebalance it against the comparatively short
storage/coupling treatment. Re-reading \texttt{Method\_and\_Implementation.md}
in full shows this was the wrong prescription: that document's own cut list
already checked subsection F against exactly this question (``what Section IV
must define is the observation $o_i(t)$ -- the induced subgraph, the potential
$\eta_i$, and the communication graph -- because those specify what an agent
is allowed to know, which is a property of the problem'') and kept it,
routing only the encoder/message-passing equations to
\ref{sec:method:model}, where \texttt{problem\_formulation\_fixed.tex}
already has them. The size imbalance documented under [G14] above is real and
numeric, but it is not evidence of misplaced Method content -- there isn't
any to move. It reflects that defining what a bounded, local observer may
know legitimately takes more apparatus than defining an abstract storage
update does; treating that as a defect and cutting it to match a line count
would itself be the kind of drift this section is supposed to resist. The
right response to the imbalance was narrower than ``trim Observations'': fix
the one place the imbalance was load-bearing, which is the congestion-signal
asymmetry above.

\subsection*{[G16] $d_G$ was defined over all of $G$, not just the occupiable subgraph}
$N^{+}(v) = \{w \in V_{\mathrm{mov}} : (v,w)\in E\}$ already restricts every
real move to land in $V_{\mathrm{mov}}$, so no agent trajectory can ever pass
through a vertex outside it. But $d_G(v,w)$ was defined as the shortest-path
cost ``along directed edges of $G$'' -- the full graph, $V_{\mathrm{mov}}$
included or not. If a future instance ever has $V_{\mathrm{mov}} \subsetneq
V$ (the environment section explicitly permits this: $V_{\mathrm{mov}}
\subseteq V$, not $=$), $d_G$ could report a shortest-path cost that routes
through a non-occupiable waypoint -- a number no real agent could ever
achieve, since every real hop is constrained to $V_{\mathrm{mov}}$. Every
downstream use of $d_G$ ($\eta_i$'s potential, the connectivity assumption,
the communication-graph radius, all of \ref{sec:pf:observations}) would then
silently be using an inadmissible distance. \textbf{Fix:} $d_G(v,w)$ is now
defined only for $v,w\in V_{\mathrm{mov}}$, over $G[V_{\mathrm{mov}}]$ (the
same induced-subgraph notation \texttt{eq:localsubgraph} already uses),
with a parenthetical noting the restriction is currently vacuous since every
generated instance so far has $V_{\mathrm{mov}}=V$. Found while explaining,
not auditing -- a reminder that ``why does this indirection exist'' questions
are worth asking even after a section has already been reviewed twice.

\subsection*{[G17] ``Free agent'' was used but never defined}
$\pi^{\mathrm{assign}}$'s own domain (\ref{sec:pf:controllers}) is stated as
``free agents and $\mathcal{Q}_t$,'' and \texttt{method\_implementation\_-
skeleton.tex}'s \texttt{eq:assignmentrule} matches ``the free agent'' to each
task -- both use the term as if self-evident, and neither states what makes
an agent free rather than occupied. Nothing else in \ref{sec:pf:tasks} names
\emph{which} agent is serving an active task either: $y_j$ (assignment
\emph{time}) was defined, but not \emph{to whom}. \textbf{Fix:} $\sigma(j)\in
\{1,\dots,m\}$ added alongside $y_j$ in \texttt{eq:lifecycle}'s paragraph
(the agent task assignment selects at $y_j$), and \texttt{eq:free} defines
free/occupied directly from $\sigma$ and $\mathcal{B}_t$: agent $a_i$ is free
at $t$ iff no active task is served by it. \texttt{eq:assignmentrule} in the
Method skeleton now points back to \texttt{eq:free} instead of relying on the
same informal term it previously left undefined.

\section{Comparison against Knippenberg, Holenderski \& Menkovski (2021)}
\label{sec:knippenberg}

\citet{knippenberg2021} (``Time-Constrained Multi-Agent Path Finding in
Non-Lattice Graphs with Deep Reinforcement Learning,'' ACML 2021) is the
closest prior-art paper to this thesis's environment and observation model:
a directed, non-lattice graph; a local observation of fixed depth $d$
centred on the agent; a per-node distance-to-goal annotation; a shared
policy over a wait-or-move action space. The thesis's problem formulation is
best read as taking that architecture and (a) embedding it inside a
lifelong/online MAPD setting with continuous task arrivals rather than
\citeauthor{knippenberg2021}'s batch/episodic one-job-per-agent setting, and
(b) coupling it to a new second layer -- storage placement -- that has no
counterpart in \citeauthor{knippenberg2021} at all. What follows is organised
by what carried over unchanged, what was correctly adapted, what is a new
extension, and one place the comparison surfaces an actual gap.

\subsection{Carried over directly (and appropriately)}
\begin{itemize}
  \item \textbf{Directed, non-lattice graph.} Identical framing in both.
  \item \textbf{Local observation of fixed depth $d$.}
    \citeauthor{knippenberg2021}'s observation is ``centered on the agent's
    current location'' with size ``determined by a global depth parameter
    $d$''; this is exactly $V_i^{(d)}(t)$, \texttt{eq:localsubgraph}.
  \item \textbf{Distance-to-goal node annotation.}
    \citeauthor{knippenberg2021} includes, per visible node, ``the distance
    of each node to the goal node,'' precomputed once per instance -- this
    is exactly $\eta_i(v,t) = d_G(v,q_i(t))$, \texttt{eq:potential}.
  \item \textbf{Wait-or-move action space, shared policy.} Identical in
    both.
\end{itemize}

\subsection{Correctly, non-trivially adapted}
\begin{itemize}
  \item \textbf{Dynamic vs.\ static target.} \citeauthor{knippenberg2021}'s
    distance-to-goal values ``need only be pre-computed once for every
    problem instance,'' because each agent has exactly one job (one fixed
    goal) for the whole episode. The thesis is lifelong: an agent's target
    changes at every reassignment, and the thesis states this explicitly
    (``$\eta_i$ itself changes whenever the target changes, which happens at
    every reassignment''). This is a correct and necessary generalisation
    the source paper didn't need to make -- a good sign, not a borrowed
    shortcut.
  \item \textbf{Vertex/swap-conflict as two orthogonal equations vs.\ one
    prose definition.} \citeauthor{knippenberg2021} defines a ``node
    collision'' (two cases in prose: moving onto an occupied node, or two
    agents moving to the same node) and an ``edge collision'' (explicitly
    qualified to \emph{bi-directional} edges only). The thesis's
    \texttt{eq:vertexconflict} covers both node-collision cases with one
    inequality, and its \texttt{eq:swapconflict} needs no ``bi-directional''
    qualifier: a swap is only physically reachable when
    \texttt{eq:transition} already requires both $(v,w)\in E$ and
    $(w,v)\in E$ to hold for the two agents involved, so the bidirectionality
    falls out of constraint composition instead of being hand-added to the
    collision definition. This is a tighter formalisation than the source.
\end{itemize}

\subsection{Genuine extensions (no counterpart in the source paper)}
\begin{itemize}
  \item \textbf{The entire storage layer} -- $x_t$, $\mathcal{X}$,
    \texttt{eq:storageupdate}, \texttt{eq:storageobjective}, and the
    coupling condition \texttt{eq:coupling} that links it back to routing.
    \citeauthor{knippenberg2021} is a single-layer routing problem; this
    second, slower-timescale layer and its feedback loop is the thesis's
    actual contribution over the source paper, and is worth stating this
    plainly wherever the thesis positions itself against MAPF/MAPD
    literature.
  \item \textbf{Explicit communication graph $G_t^{\mathcal{A}}$ and
    messages.} \citeauthor{knippenberg2021} states explicitly they consider
    only ``implicit communication'' (agents observe each other's state, but
    nothing is explicitly communicated). \texttt{eq:commgraph} and the
    message component of \texttt{eq:observation} are a genuine addition.
  \item \textbf{The congestion feature $\delta_i(t)$.} No counterpart in
    the source.
  \item \textbf{Three information regimes (centralised / section-based /
    decentralised) with a shared assignment mechanism.}
    \citeauthor{knippenberg2021} studies a single decentralised regime and
    has no assignment sub-problem at all (each agent's one job is given, not
    chosen online) -- the thesis's Table~I and the assignment/routing split
    are new, driven by the lifelong setting's need to keep deciding who does
    what.
  \item \textbf{Throughput, backlog, entropy/concentration.} All specific
    to the continuous/lifelong framing; \citeauthor{knippenberg2021}'s
    objective (a single scalar per solved batch instance, eq.~1--2 in the
    source) has no analogue for an ongoing system with no end state.
\end{itemize}

\subsection{[K1, Deviation] Strong connectivity}
\citeauthor{knippenberg2021} state plainly: ``environment graphs are assumed
to be strongly connected, i.e.\ each node is reachable from all other
nodes.'' The thesis (even after [G10]'s fix) states the weaker ``$d_G(v,w)$
is finite for every pair of vertices that appears in the constructions
below.'' The source's version is simpler to state and to verify computationally
(one graph-level check, independent of what gets used later); the thesis's
version is strictly weaker and therefore more permissive of the ``irregular
blocks'' / unreachable side-room layouts its own motivating paragraph gestures
at. \textbf{This is a judgement call, not a correctness issue}: if generated
instances never actually exercise a disconnected pocket, the weaker statement
buys nothing and the source's crisper assumption should probably be adopted
outright; if partial reachability is a deliberate stress case for the
decentralised controller (an agent whose target briefly becomes unreachable),
the weaker version is doing real work and is worth keeping, explicitly framed
as a deviation from the more common assumption. The fixed .tex takes the
conservative option -- keeps the weaker statement, adds one sentence
contrasting it with \citeauthor{knippenberg2021}'s choice and naming the
tradeoff -- and leaves the actual decision to Section~VI (\emph{Instance
generation}), where it is checkable against what instances are actually
generated.

\subsection{[K2, Gap] No task deadlines}
This is the most substantial finding from the comparison. Time constraints
are \citeauthor{knippenberg2021}'s entire contribution -- their title is
literally ``Time-\emph{Constrained} MAPF'' -- realised as a job 4-tuple
$(t_i^s, t_i^d, v_i^s, v_i^g)$ with an explicit deadline $t_i^d$, and an
objective that is either a hard feasibility cut (a plan that misses its
deadline is invalid, eq.~1 in the source) or a soft penalty term
$w\cdot\max(0, t_{\pi,\text{actual}} - t_{\pi,\text{deadline}})$ added to the
cost (eq.~2). The thesis's task tuple \texttt{eq:task} carries only a release
time $r_j$; there is no deadline anywhere, and timeliness is scored only in
aggregate, after the fact, via the mean $\bar\zeta_T$ (\texttt{eq:throughput}).
Under the thesis's formulation, no individual task can be said to have
\emph{failed} -- an arbitrarily late delivery only nudges an average.

Given the thesis explicitly reuses \citeauthor{knippenberg2021}'s
architecture (built specifically to handle deadlines) while dropping the one
thing that architecture was built for, this needed to be an explicit,
stated choice rather than a silent omission -- warehouse SLAs are a real
business constraint, so a reader familiar with the source paper would
reasonably ask why deadlines didn't carry over too. \textbf{Fix (already
applied in \texttt{problem\_formulation\_fixed.tex}, Section~J):} one
paragraph stating the omission explicitly, alongside the other scope
exclusions, and naming what reintroducing deadlines would require: a
$t_j^{\mathrm{d}}$ field added to \texttt{eq:task}, plus either a feasibility
condition (hard) or a term in \texttt{eq:functional} of the same shape as
\citeauthor{knippenberg2021}'s eq.~2 (soft). Both are Problem-Formulation-level
changes, not Method-level ones, since either would redefine what a valid or
well-scored solution \emph{is} -- consistent with how the thesis already
treats \texttt{eq:coupling} and the conflict constraints as problem-level
because they define validity, not merely score it.

\subsection{[K3--K6, Deviation] Smaller, favourable deviations}
\begin{itemize}
  \item \textbf{[K3] Edge cost $c(e) > 0$ (thesis) vs.\ $c(v,v') \ge 0$
    (source).} The thesis's strict positivity rules out zero-cost cycles,
    which keeps $d_G$ well-behaved as a strictly hop-sensitive quantity and
    keeps the corrected potential-bound argument ([Error] above) free of
    degenerate edge cases. A defensible tightening, not a regression.
  \item \textbf{[K4] No \texttt{canWait}-style per-node restriction.}
    \citeauthor{knippenberg2021} give each node a boolean attribute
    controlling whether waiting is allowed there (e.g.\ a mainline track
    where a train cannot simply stop); the thesis assumes waiting is legal
    everywhere via a single $c_{\mathrm{wait}}$. This underwrites the
    thesis's ``a collision-free joint action always exists'' claim
    (Section~B) for free, but is a real loss of expressiveness relative to
    the source if a future revision wants to model genuinely
    non-stoppable segments -- worth a one-line note in Scope if that
    generality is ever wanted; not fixed here since it is not currently
    claimed as a capability.
  \item \textbf{[K5] No dummy-node start-time mechanism needed.}
    \citeauthor{knippenberg2021} materialise each agent at a private dummy
    node until its individual start time $t_i^s$ elapses, to avoid an agent
    appearing mid-graph and causing a spurious collision. The thesis avoids
    this entirely by construction: the fleet $\mathcal{A}$ is fixed and
    persistent (agents idle at $V_{ep}$ between tasks), and it is
    \emph{tasks}, not agents, that have release times $r_j$. A structural
    consequence of the lifelong framing, not a borrowed mechanism -- correctly
    absent, not simply omitted.
  \item \textbf{[K6] Reward design (Method-level, noted here for
    completeness).} \citeauthor{knippenberg2021} use a fixed table of
    empirically-tuned per-action penalties (Table~1 in the source: e.g.\
    move $-0.1$, collision $-2$, goal completion $+10$) with no stated
    invariance guarantee. \texttt{Method\_and\_Implementation.md} instead
    specifies potential-based shaping, $\gamma\Phi_i(t{+}1)-\Phi_i(t)$ with
    $\Phi_i(t) = -\eta_i(\ell_i(t),t)$, which provably preserves the
    underlying game's equilibria (\citealp{ng1999policy}; extended to
    dynamic potentials by \citealp{devlin2012dynamic}) -- directly relevant
    here since $\Phi_i$ is non-stationary by construction (the target $q_i(t)$
    changes at every reassignment). \citeauthor{knippenberg2021} even flag
    their own use of DQN as a limitation and name PPO as future work (their
    Section~8); \texttt{Method\_and\_Implementation.md}'s planned PPO
    training procedure directly answers that. Correctly scoped to Method, not
    Problem Formulation -- noted here only because the comparison is what
    surfaced it as a deliberate, citable improvement rather than an
    arbitrary choice.
\end{itemize}

\subsection{[K7, Clarity] Cost-asymmetry sentence was untethered, and $c(v,w)$ was never tied to $c(e)$}
An earlier revision of \texttt{problem\_formulation\_fixed.tex} added a
standalone sentence under the \texttt{[G10]} tag stating that
$(v,w)\in E$ does not imply $(w,v)\in E$ and that the two directions ``may
carry different costs.'' Directional cost asymmetry is not wrong --
\citeauthor{knippenberg2021}'s own formalism defines cost per ordered pair
$c(v,v')$ on a directed graph too -- but it is a free corollary of
directedness, not an independent modelling choice, and no generated instance
in either this thesis or \citeauthor{knippenberg2021} actually exercises a
bidirectional edge with unequal costs in each direction (their bidirectional
edges are drawn as a single shared cost, e.g.\ Figure~1's ``cost~$=1$''). The
sentence was therefore carrying a \texttt{[G10]} tag without matching
\texttt{[G10]}'s own rationale in this document, i.e.\ an undocumented claim.
\textbf{Fix:} folded into the directedness explanation itself (``directedness
matters concretely here \ldots aisles are often one-way, so $(v,w)\in E$ does
not by itself imply $(w,v)\in E$, nor does it imply the two directions cost
the same where both exist''), so the asymmetry reads as a consequence of
being directed rather than a separately asserted feature.

Separately, $c(v,w)$ was used three times (in \texttt{eq:onestepcost} and
\texttt{eq:potentialbound}) without ever being identified with $c(e)$ for
$e=(v,w)$, the only place the cost function is actually defined -- a gap the
section's own notation-disambiguation paragraph (\texttt{sec:pf:notation})
did not catch, since it lists $c(e)$, $\hat c(v,w)$, $\bar c_T$, $C_T$ and
$\mathcal{C}_t$ as the family needing disambiguation but never states that
plain $c(v,w)$ \emph{is} $c(e)$. \textbf{Fix:} one clause added where $c(e)$
is first defined -- ``written $c(v,w)$ when $e=(v,w)$'' -- closing the gap
without introducing a sixth symbol into the disambiguation table.

\subsection*{[K8, Deviation] Action masking vs.\ Knippenberg's penalised-invalid-move handling}
Re-reading \S3 of \citeauthor{knippenberg2021} closely: ``If an agent chooses
a move action that is not valid, i.e.\ a move action outside the range of
outgoing edges of its current location, the action is treated as a `wait'
action, and it is penalised accordingly.'' Their fixed-size action head can
select a nonexistent edge; the environment silently reinterprets that choice
as a wait and lets the reward signal teach the policy not to do that. This
thesis's \texttt{eq:mask} (\texttt{method\_implementation\_skeleton.tex},
\ref{sec:method:rl}) instead sets the logits for nonexistent moves to
$-\infty$ before the softmax, so an illegal move has probability zero rather
than a penalised nonzero one -- legal \emph{by construction}, not learned by
penalty. This was already the right design and already correctly cross-tied
to $\mathcal{U}(v)$ vs.\ $\mathcal{U}_i$ in \ref{sec:pf:game} (the
\texttt{[Gap]} entry above), but the comparison to
\citeauthor{knippenberg2021}'s own mechanism -- and the actual argument for
why masking is preferable, not just different -- lived only in this
document, not in the thesis text itself. \textbf{Fix:} two sentences added
to the masking paragraph naming \citet{knippenberg2021}'s penalised
approach explicitly and stating why masking is preferred: it never spends
training signal on moves already known illegal from $N^{+}(\ell_i(t))$, it
keeps the executed action identical to the sampled one (their reinterpretation
lets a sampled ``move'' execute as ``wait,'' which the value function then
has to reconcile), and it costs nothing extra since $N^{+}(v)$ is already
required to define $\mathcal{U}(v)$ regardless. Same lesson as \texttt{[K6]}:
a citable improvement over the source is only citable if the thesis actually
says so.

\section{Comparison against Ma, Li, Kumar \& Koenig (2017)}
\label{sec:ma2017}

\citet{ma2017lifelong} (``Lifelong Multi-Agent Path Finding for Online
Pickup and Delivery Tasks,'' AAMAS 2017) is where the term MAPD and the
notion of a well-formed MAPD instance both come from; this thesis uses both
throughout without, until now, ever citing the source or fully carrying over
what it defines.

\subsection*{[M1, Gap] MAPD and well-formedness cited but not attributed}
Two uncredited borrowings: ``Multi-Agent Pickup and Delivery (MAPD)'' at its
first use in the Introduction, and the well-formed-instances condition in
\ref{sec:pf:env} (at least $m$ non-task endpoints; no path between two
endpoints blocked by a third), which is \citeauthor{ma2017lifelong}'s
Definition~1, restated in this thesis's own words but previously with no
citation at all. \textbf{Fix:} \citep{ma2017lifelong} added at MAPD's first
use; the well-formedness paragraph now reads ``\citet{ma2017lifelong}'s
well-formed MAPD instances'' rather than stating the condition as though it
were invented here. A new three-panel figure, \texttt{fig:wellformed},
mirrors \citeauthor{ma2017lifelong}'s own Figure~2 (which contrasts a
well-formed instance against two ways the condition can fail) rather than
inventing a new pedagogical device: (a) well-formed, (b) too few non-task
endpoints, (c) a path between two endpoints blocked by a third.

\subsection*{[M1, continued] Re-tasking: permitted in the source, silently absent here}
\citeauthor{ma2017lifelong} explicitly allow an agent to be re-tasked --
reassigned to a different job -- any time before it reaches its current
task's pickup location, and give the reason: ``agents can often be re-tasked
before they have picked up a good but have to deliver it afterward.'' This
thesis's \texttt{eq:lifecycle} moves a task from $\mathcal{Q}_t$ to
$\mathcal{B}_t$ at $y_j$ with no path back, which \emph{implies} no
re-tasking, but nothing said so before now. Given the source paper is cited
specifically for the MAPD problem class, a reader who knows it would
reasonably assume re-tasking carried over too. \textbf{Fix:} one paragraph
added to \ref{sec:pf:scope} stating the narrowing explicitly, using the
$\sigma$ notation from [G17]: $\sigma(j)$ is fixed for the lifetime of
$\mathcal{B}_t$-membership, so re-tasking is a Problem-Formulation-level
change (it would let $\pi^{\mathrm{assign}}$ revise $\sigma$ for
already-assigned, not-yet-picked-up tasks), not a Method-level one --
consistent with how deadlines ([K2]) are already treated as the same kind of
change.

\subsection*{[M2, Gap] The restated well-formedness clause checked a narrower endpoint set than the source's}
\citeauthor{ma2017lifelong}'s own $V_{\mathrm{ep}}$ already contains every
pickup and delivery location ($V_{\mathrm{tsk}} \subseteq V_{\mathrm{ep}}$),
so their path-condition clause (c) -- ``for any two endpoints, a path exists
that traverses no other endpoints'' -- ranges over task \emph{and} non-task
endpoints together. This thesis keeps $V_{\mathrm{str}}$, $V_{\mathrm{del}}$
and $V_{\mathrm{ep}}$ as three separate, typically-disjoint roles (a
deliberate and more general choice than the source's nesting, discussed on
its own merits elsewhere), which means the restated condition -- ``between
any pair of endpoints some path avoids every third endpoint'' -- silently
ranged over a narrower set than \citeauthor{ma2017lifelong}'s original if
``endpoint'' were read as this section's $V_{\mathrm{ep}}$ alone: it said
nothing about a path between, say, two storage vertices, even though the
source's completeness guarantee needs exactly that. \textbf{Fix:} one clause
added stating that for this specific check, ``endpoint'' means
$V_{\mathrm{ep}} \cup V_{\mathrm{str}} \cup V_{\mathrm{del}}$ together,
matching \citeauthor{ma2017lifelong}'s convention for exactly the well-
formedness test and nowhere else -- $V_{\mathrm{str}}$, $V_{\mathrm{del}}$
and $V_{\mathrm{ep}}$ remain separate roles everywhere else in the section,
so this does not quietly redefine any of the three. Found by working through
what ``more general than the source, not more limited'' actually requires
of a condition borrowed \emph{from} the source: generality in the role
structure does not automatically carry over to a check stated in terms of
those roles unless it is checked.

\subsection*{[M3, Gap] ``A path exists'' needs a direction on a directed graph}
\citeauthor{ma2017lifelong}'s graph is undirected, so their clause (c) --
``a path exists between any two endpoints'' -- is automatically symmetric: a
path from $v$ to $w$ \emph{is} a path from $w$ to $v$, walked backwards. This
thesis's $G$ is directed (\ref{sec:pf:env}, and the whole point of the
[G10]/[K1] departure from Knippenberg's own graph choice), so the same
sentence is ambiguous here: does the check need a path from $v$ to $w$, from
$w$ to $v$, or both? Since Token Passing-style algorithms route an agent
between whichever pair of endpoints a task assignment happens to produce --
not a fixed direction chosen in advance -- only requiring \emph{both}
directions preserves the completeness guarantee the citation is for; a
single-direction reading would silently be a weaker (and wrong) condition.
\textbf{Fix:} the path-condition clause now states explicitly that, because
$G$ is directed, it requires a path from $v$ to $w$ \emph{and} a path from
$w$ to $v$ for every pair, each avoiding every third endpoint. Found in the
same conversation as [M2] and for the same underlying reason: a property
proved on the source's graph does not automatically transfer to a strictly
more general graph class without checking which of the source's implicit
assumptions (here, undirectedness) the property actually used.

\subsection*{[M4, Deviation] The active phase starts at assignment, not pickup arrival, and this thesis never said so}
\citeauthor{ma2017lifelong} state plainly: ``When the agent reaches the
pickup location, it starts to execute the task and the task is removed from
$\mathcal{T}$'' -- their task set only drops a task, i.e.\ their active
phase only begins, at \emph{pickup arrival}, not at assignment. This
thesis's \texttt{eq:lifecycle} already uses a different boundary --
$\mathcal{B}_t$ begins at $y_j$ (assignment) -- and the surrounding prose
already explicitly rejects ``arrival of the agent at $s_j$'' as the
boundary. What it did not do, before now, is note that the rejected
alternative is \citeauthor{ma2017lifelong}'s own convention, or say why this
thesis picks the earlier one. Given the source is now cited repeatedly for
this exact lifecycle structure, a reader who knows it would notice the
silent divergence. \textbf{Fix:} one clause naming
\citet{ma2017lifelong}'s convention and the reason for departing from it:
using $y_j$ rather than pickup arrival means time lost travelling to $s_j$
is attributed to the active phase and is therefore eligible for the
blocking measure $W_T$ \eqref{eq:waiting}, instead of being folded into
generic pre-assignment queueing (which $B_T$'s $|\mathcal{Q}_T|$ term
already covers). This is a considered choice, not an oversight -- but it
was an unstated one until this pass.

\subsection*{[M4, continued] This same boundary is why \texttt{[G17]} and \texttt{[M1]} agree with each other, not just with the source}
Re-checking the section once more surfaces that \texttt{[G17]}'s
\texttt{eq:free} is built directly on $\mathcal{B}_t$, which this entry just
pinned to $y_j$. So an agent becomes \emph{occupied} the instant it is
assigned, not when it reaches the pickup -- and that is precisely why
\texttt{[M1]}'s re-tasking prohibition holds \emph{before} pickup as well as
after: an assigned-but-travelling agent is already occupied under this
thesis's boundary, so there is no ``free'' window in which
\citeauthor{ma2017lifelong}'s reassignment could happen even if it were
permitted. In their model, the later (pickup-arrival) boundary is exactly
what leaves that window open, which is why they can allow it and this
thesis's own choice of boundary is why it cannot. \texttt{[G17]},
\texttt{[M1]} and \texttt{[M4]} are not three independent findings that
happen to be consistent -- they are one decision (where $\mathcal{B}_t$
begins) examined from the free/occupied side, the re-tasking side, and the
$W_T$-measurement side respectively. No further fix: the three passages are
already individually correct and mutually consistent; this is a note for
whoever reads GAPS.tex next; so the connection does not have to be
rediscovered from scratch.

\section{Where does this leave Method\_and\_Implementation.md?}
\label{sec:method-mapping}

Checking \texttt{problem\_formulation.tex} (as given, not a hypothetical
denser draft) against \texttt{Method\_and\_Implementation.md}'s own
six-item cut list:

\begin{enumerate}
  \item \emph{Controller-architecture algorithmic detail (prioritised
    planning, matching costs, zone hand-off protocols)} -- \textbf{not
    present.} Section~E as given only states the abstract information
    signature and Table~I; no algorithm-level content to cut.
  \item \emph{Encoder / message-passing equations
    ($z_i=\mathrm{Enc}_G(\cdot)$, $\phi_q,\psi_q$)} -- \textbf{not present.}
    Section~F as given defines only $V_i^{(d)}$, $\eta_i$, $\delta_i$,
    $G_t^{\mathcal{A}}$, $o_i(t)$ -- no architecture equations.
  \item \emph{Exact functional forms of $D$, $K$, $R$ and weights
    $\alpha,\beta,\chi$} -- \textbf{already deferred} in the text as given
    (``specified in Section~\ref{sec:method:storage}'').
  \item \emph{The reward function} -- \textbf{already deferred}
    (\texttt{eq:objective}'s prose: ``$R_i$ are left unspecified here'').
  \item \emph{Defensive justification passages} -- \textbf{present, low
    priority.} These exist (paraphrased in
    \texttt{Method\_and\_Implementation.md} rather than quoted verbatim) but
    are one-line footnote-style asides, not misplaced algorithmic content;
    rephrasing them is a copy-edit, not a structural fix, and wasn't applied
    here.
  \item \emph{Full notation table $\to$ appendix} -- \textbf{applies, but
    only to the table added in this review} (the original had none). Fixed:
    see [G11] above and Appendix~\ref{app:notation} below.
\end{enumerate}

The short version: \textbf{almost nothing from this review belongs in
Method\_and\_Implementation.md.} Every fix in
\texttt{problem\_formulation\_fixed.tex} is definitional (``what the object
is''), which is exactly what Section~IV is supposed to contain per that
document's own contract table. The one piece of content that needed to move
doesn't belong in Method either -- by \texttt{Method\_and\_Implementation.md}'s
own instruction, a full notation table is \emph{Appendix} material, a third
location distinct from both Problem Formulation and Method. That relocation
is done in Appendix~\ref{app:notation} of this document (ready to paste into
the thesis's own appendix chapter, which is not a file present in this
repository).

\appendix
\section{Notation}
\label{app:notation}

Two tables. The first is the full glossary requested during review: every
symbol introduced in Section~IV, in the order it first appears, with a
plain-language gloss and where it is defined. The second (unchanged from the
earlier version of this appendix) is the narrower disambiguation table for
symbol families that share a base letter across unrelated objects -- keep
both, since the first is for looking a symbol up and the second is for the
specific worry ``is this the same $c$ as that $c$.''

\begin{table}[h]
  \centering
  \caption{Every symbol used in Section~IV, in order of first appearance.}
  \label{tab:notation-full}
  \begin{tabular}{ll}
    \toprule
    Symbol & Meaning (first defined at) \\
    \midrule
    \multicolumn{2}{l}{\textit{The environment (\ref{sec:pf:env})}} \\
    $G=(V,E)$ & the directed warehouse graph \\
    $V_{\mathrm{mov}}$ & vertices an agent may occupy \\
    $V_{\mathrm{str}}$ & storage faces (pickup locations) \\
    $V_{\mathrm{del}}$ & delivery points \\
    $V_{\mathrm{ep}}$ & endpoints: safe to idle at without blocking transit \\
    $c(e)=c(v,w)$ & cost of traversing edge $e=(v,w)$, eq.~(env def.) \\
    $d_G(v,w)$ & cheapest directed path cost, over $G[V_{\mathrm{mov}}]$ only \\
    $N^{+}(v)$ & movable successors of $v$ \\
    $\mathcal{U}(v)$ & legal actions at vertex $v$, eq.~(actions) \\
    $c_{\mathrm{wait}}$ & cost of waiting one timestep \\
    $\hat c(v,w)$ & realised one-step cost (handles waiting too), eq.~(onestepcost) \\
    $m=|\mathcal{A}|$ & fleet size (defined in \ref{sec:pf:agents}, used here) \\
    \addlinespace
    \multicolumn{2}{l}{\textit{Agents and collisions (\ref{sec:pf:agents})}} \\
    $\mathcal{A}=\{a_1,\dots,a_m\}$ & the fleet \\
    $\ell_i(t)$ & agent $a_i$'s location at time $t$ \\
    $u_i(t)$ & agent $a_i$'s chosen action at time $t$ \\
    $p_i$ & agent $a_i$'s realised route over the horizon \\
    \addlinespace
    \multicolumn{2}{l}{\textit{Where the items are (\ref{sec:pf:storage})}} \\
    $\mathcal{K}$ & set of SKU types \\
    $x_t$ & storage configuration: units of each SKU at each storage vertex \\
    $b_k$ & shelf capacity one unit of SKU $k$ consumes \\
    $\mathrm{cap}(v)$ & shelf capacity of storage vertex $v$ \\
    $\mathcal{X}$ & feasible storage configurations, eq.~(feasiblestorage) \\
    $x_0$ & initial storage configuration, supplied with the instance \\
    $F$ & the storage update rule \\
    $\Delta$ & storage epoch length (the slow clock's period) \\
    $\hat\rho_t,\hat\mu_t,\hat w_t$ & causal estimates: demand, edge traffic, waiting \\
    $F_{\mathrm{fix}}$ & the do-nothing baseline, $F_{\mathrm{fix}}(x,\cdot)=x$ \\
    $\mathcal{F}$ & the set of admissible storage rules \\
    $D,K,R$ & storage objective terms: access distance, congestion, reassignment \\
    $\alpha,\beta,\chi$ & weights on $D,K,R$ in eq.~(storageobjective) \\
    \addlinespace
    \multicolumn{2}{l}{\textit{Where the work comes from (\ref{sec:pf:tasks})}} \\
    $P_{\mathrm{ord}}$ & the order-arrival distribution \\
    $\tau_j=(r_j,s_j,g_j,k_j)$ & a task: release time, pickup, delivery, requested SKU \\
    $\lambda_{\mathrm{task}}$ & expected tasks released per timestep \\
    $y_j$ & assignment time of $\tau_j$ \\
    $\sigma(j)$ & the agent assigned to $\tau_j$ at $y_j$ \\
    $d_j$ & completion (delivery) time of $\tau_j$ \\
    $\mathcal{Q}_t,\mathcal{B}_t,\mathcal{C}_t$ & waiting / active / completed task sets, eq.~(lifecycle) \\
    $\zeta_j$ & service time of $\tau_j$, eq.~(servicetime) \\
    \addlinespace
    \multicolumn{2}{l}{\textit{Who decides (\ref{sec:pf:controllers})}} \\
    $\pi=(\pi^{\mathrm{assign}},\pi^{\mathrm{route}})$ & a controller: assignment map + routing map \\
    $\Pi$ & the set of admissible controllers \\
    $\mathcal{Y}=\{Y_1,\dots,Y_p\}$ & the zone partition for section-based control \\
    $\delta(S,t;\mathrm{excl})$ & occupancy fraction of window $S$, excluding agents in $\mathrm{excl}$ \\
    $\theta$ & a learned controller's parameters \\
    \addlinespace
    \multicolumn{2}{l}{\textit{What a decentralised agent sees (\ref{sec:pf:observations})}} \\
    $d$ & observation depth \\
    $V_i^{(d)}(t),\mathcal{G}_i^{(d)}(t)$ & visible vertices / induced subgraph, eq.~(localsubgraph) \\
    $q_i(t)$ & agent $a_i$'s current target vertex \\
    $\eta_i(v,t)$ & potential: distance from $v$ to $q_i(t)$, eq.~(potential) \\
    $r_{\mathrm{cng}}$ & congestion radius \\
    $\delta_i(t)$ & decentralised instance of $\delta(S,t;\mathrm{excl})$, eq.~(congestion) \\
    $r_{\mathrm{com}}$ & communication radius \\
    $G^{\mathcal{A}}_t=(\mathcal{A},E^{\mathcal{A}}_t)$ & who-can-talk-to-whom graph, eq.~(commgraph) \\
    $o_i(t)$ & the full observation agent $a_i$ receives, eq.~(observation) \\
    \addlinespace
    \multicolumn{2}{l}{\textit{The learning problem (\ref{sec:pf:game})}} \\
    $\mathcal{M}$ & the POSG tuple, eq.~(posg) \\
    $\mathcal{S}$ & the state space \\
    $P$ & the transition kernel \\
    $\mathcal{O}_i,O_i$ & observation space / observation function \\
    $R_i$ & agent $a_i$'s reward function \\
    $\gamma$ & discount factor \\
    $\mathcal{U}_i$ & agent $a_i$'s fixed action space in $\mathcal{M}$ (union over vertices) \\
    $J(\theta)$ & the team-level learning objective, eq.~(objective) \\
    \addlinespace
    \multicolumn{2}{l}{\textit{How the two layers close the loop (\ref{sec:pf:coupling})}} \\
    $s_t$ & the full system state (agents, task sets, storage config.) \\
    \addlinespace
    \multicolumn{2}{l}{\textit{How a run is scored (\ref{sec:pf:measures})}} \\
    $\Lambda_T$ & throughput, eq.~(throughput) \\
    $\bar\zeta_T$ & mean service time \\
    $\bar c_T$ & mean movement cost per completed task, eq.~(movementcost) \\
    $\mu_T(e)$ & traversal count of edge $e$ over the evaluation window \\
    $E^{+}_T$ & edges with at least one traversal \\
    $p_T(e)$ & normalised traversal share of edge $e$ \\
    $H_T$ & normalised entropy of edge usage, eq.~(entropy) \\
    $C_T$ & concentration index, $1-H_T$, eq.~(concentration) \\
    $\kappa_T$ & mean per-decision controller runtime, eq.~(runtime) \\
    $B_T$ & backlog, $|\mathcal{Q}_T|+|\mathcal{B}_T|$ \\
    $W_T$ & blocking measure (time lost after assignment), eq.~(waiting) \\
    $\omega_j(t)$ & indicator that the agent serving $\tau_j$ is blocked at $t$ \\
    $\mathcal{J}(F,\pi)$ & the system-level criterion, eq.~(functional) \\
    $q_{\mathrm{svc}},q_{\mathrm{mov}},q_{\mathrm{wait}},q_{\mathrm{conc}},q_{\mathrm{back}}$ & weights in $\mathcal{J}$ \\
    \bottomrule
  \end{tabular}
\end{table}

\noindent As with the disambiguation table below, replace each
\texttt{eq.~(xxx)} with \verb|\eqref{eq:xxx}| when pasting this into the
thesis's own appendix -- labels already match
\texttt{problem\_formulation\_fixed.tex} verbatim.

Several symbol families in Section~IV share a base letter across unrelated
objects. None are typos or the same object under different names; each pair
below denotes two genuinely different things that happen to look alike.

\begin{table}[h]
  \centering
  \caption{Symbols that share a base letter but denote unrelated objects.}
  \label{tab:notation}
  \begin{tabular}{ll}
    \toprule
    Symbols & Distinct meanings \\
    \midrule
    $c(e)$, $c_{\mathrm{wait}}$ & edge cost / wait cost (instance parameters) \\
    $\hat{c}(v,w)$ & realised one-step cost, eq.~(onestepcost) \\
    $\bar{c}_T$ & mean movement cost per completed task, eq.~(movementcost) \\
    $C_T$ & concentration index, eq.~(concentration) \\
    $\mathcal{C}_t$ & completed-task set, eq.~(lifecycle) \\
    \addlinespace
    $\mathcal{Q}_t$ & waiting-task set, eq.~(lifecycle) \\
    $q_i(t)$ & agent $i$'s current target vertex, eq.~(potential) \\
    $q_{\mathrm{svc}}, q_{\mathrm{mov}}, q_{\mathrm{wait}}, q_{\mathrm{conc}}, q_{\mathrm{back}}$
      & weights in the system-level criterion, eq.~(functional) \\
    \addlinespace
    $\mathcal{O}_i$ & observation space, eq.~(posg) \\
    $O_i$ & observation function, eq.~(posg) \\
    \addlinespace
    $\mathcal{U}(v)$ & action set at vertex $v$, eq.~(actions) \\
    $\mathcal{U}_i = \bigcup_{v} \mathcal{U}(v)$ & agent $i$'s fixed action space in $\mathcal{M}$, eq.~(posg) \\
    \bottomrule
  \end{tabular}
\end{table}

\noindent When pasted into the thesis's actual appendix, replace each
\texttt{eq.~(xxx)} above with \verb|\eqref{eq:xxx}| -- the labels already
match \texttt{problem\_formulation\_fixed.tex} verbatim, so this table can be
dropped in without renumbering anything.

\begin{thebibliography}{9}
\bibitem[van Knippenberg et~al.(2021)]{knippenberg2021}
Marijn van Knippenberg, Mike Holenderski, and Vlado Menkovski.
\newblock Time-Constrained Multi-Agent Path Finding in Non-Lattice Graphs
  with Deep Reinforcement Learning.
\newblock In \emph{Proceedings of the 13th Asian Conference on Machine
  Learning}, volume 157 of \emph{PMLR}, 2021.

\bibitem[Ng et~al.(1999)]{ng1999policy}
Andrew~Y. Ng, Daishi Harada, and Stuart Russell.
\newblock Policy Invariance Under Reward Transformations: Theory and
  Application to Reward Shaping.
\newblock In \emph{ICML}, 1999.

\bibitem[Devlin and Kudenko(2012)]{devlin2012dynamic}
Sam~Michael Devlin and Daniel Kudenko.
\newblock Dynamic Potential-Based Reward Shaping.
\newblock In \emph{AAMAS}, 2012.

\bibitem[Ma et~al.(2017)]{ma2017lifelong}
Hang Ma, Jiaoyang Li, T.~K. Satish Kumar, and Sven Koenig.
\newblock Lifelong Multi-Agent Path Finding for Online Pickup and Delivery
  Tasks.
\newblock In \emph{Proceedings of the 16th International Conference on
  Autonomous Agents and Multiagent Systems (AAMAS)}, pages 837--845, 2017.
\end{thebibliography}

\end{document}
